\begin{thebibliography}{999}\footnotesize
\bibitem{ref1} GBD 2023 IPV and SVAC Collaborators. Disease burden attributable to intimate partner violence against females and sexual violence against children in 204 countries and territories, 1990-2023: a systematic analysis for the Global Burden of Disease Study 2023. Lancet (London, England). 2026. PMID: 41386261.
\bibitem{ref2} Sardinha L. Global, regional, and national prevalence estimates of physical or sexual, or both, intimate partner violence against women in 2018. Lancet (London, England). 2022. PMID: 35182472.
\bibitem{ref3} Kuswanto H. Prevalence of and factors associated with female child marriage in Indonesia. PLoS ONE. 2024. PMID: 38968277.
\bibitem{ref4} Rumble L. An empirical exploration of female child marriage determinants in Indonesia. BMC Public Health. 2018. PMID: 29587705.
\bibitem{ref5} Islam MK. Attitudes to and experiences of intimate partner violence among Rohingya women who married before eighteen years of age. Global Health Action. 2021. PMID: 34323166.
\bibitem{ref6} Hawkins AJ. Does marriage and relationship education work? A meta-analytic study. Journal of Consulting and Clinical Psychology. 2008. PMID: 18837590.
\bibitem{ref7} Hawkins AJ. How effective are ACF-funded couple relationship education programs? A meta-analytic study. Family Process. 2022. PMID: 35040124.
\bibitem{ref8} Karantzas Gery C. Do relationship education programs reduce relationship aggression? A meta-analytic study. Clinical Psychology Review. 2023. PMID: 37499317.
\bibitem{ref9} Markman HJ. A randomized clinical trial of the effectiveness of premarital intervention: moderators of divorce outcomes. Journal of Family Psychology. 2013. PMID: 23421844.
\bibitem{ref10} Williamson HC. Risk moderates the outcome of relationship education: a randomized controlled trial. Journal of Consulting and Clinical Psychology. 2015. PMID: 25664643.
\bibitem{ref11} Benedetto Chiara. FIGO Preconception Checklist: Preconception care for mother and baby. International Journal of Gynaecology and Obstetrics. 2024. PMID: 38426290.
\bibitem{ref12} Dilli Priscilla Peter. Update on the practice of premarital screening for sickle cell traits in Africa: a systematic review and meta-analysis. BMC Public Health. 2024. PMID: 38822327.
\bibitem{ref13} Anwar. Characteristics and Effectiveness of Premarital Thalassaemia Educational Interventions Worldwide: A scoping review. Sultan Qaboos University Medical Journal. 2025. PMCID: PMC12445309; not PubMed-indexed.
\bibitem{ref14} Stacey Dawn. Decision aids for people facing health treatment or screening decisions. Cochrane Database of Systematic Reviews. 2024. PMID: 38284415.
\bibitem{ref15} Elwyn G. A three-talk model for shared decision making: multistage consultation process. BMJ. 2017. PMID: 29109079.
\bibitem{ref16} Upadhyay Ushma D. Development and validation of a reproductive autonomy scale. Studies in Family Planning. 2014. PMID: 24615573.
\bibitem{ref17} McCauley Heather L. Psychometric properties and refinement of the Reproductive Coercion Scale. Contraception. 2017. PMID: 27639927.
\bibitem{ref18} Upadhyay Ushma D. Development and Validation of the Sexual and Reproductive Empowerment Scale for Adolescents and Young Adults. The Journal of Adolescent Health. 2021. PMID: 32690468.
\bibitem{ref19} Dagher Myriam. Adaptation and psychometric assessment of a sexual and reproductive empowerment scale in Arabic among refugee and non-refugee adolescent girls. BMC Medical Research Methodology. 2024. PMID: 39266993.
\bibitem{ref20} Alomair N. Factors influencing sexual and reproductive health of Muslim women: a systematic review. Reproductive Health. 2020. PMID: 32138744.
\bibitem{ref21} Spencer CN. Health effects associated with exposure to intimate partner violence against women and childhood sexual abuse: a burden of proof study. Nature Medicine. 2023. PMID: 38081957.
\bibitem{ref22} Oram S. The Lancet Psychiatry Commission on intimate partner violence and mental health: advancing mental health services, research, and policy. The Lancet Psychiatry. 2022. PMID: 35569504.
\bibitem{ref23} Bacchus LJ. Interventions that prevent or respond to intimate partner violence against women and violence against children: a systematic review. The Lancet Public Health. 2024. PMID: 38702097.
\bibitem{ref24} Abramsky T. Findings from the SASA! Study: a cluster randomized controlled trial to assess the impact of a community mobilization intervention to prevent violence against women and reduce HIV risk in Kampala, Uganda. BMC medicine. 2014. PMID: 25248996.
\bibitem{ref25} Abramsky T. Ecological pathways to prevention: How does the SASA! community mobilisation model work to prevent physical intimate partner violence against women? BMC public health. 2016. PMID: 27084116.
\bibitem{ref26} Dunkle K. Effective prevention of intimate partner violence through couples training: a randomised controlled trial of Indashyikirwa in Rwanda. BMJ global health. 2020. PMID: 33355268.
\bibitem{ref27} Shaw B. Shifting Norms in Faith Communities to Reduce Intimate Partner Violence: Results from a Cluster Randomized Controlled Trial in Nigeria. Journal of interpersonal violence. 2023. PMID: 37329160.
\bibitem{ref28} Chatterji S. Community activism as a strategy to reduce intimate partner violence (IPV) in rural Rwanda: Results of a community randomised trial. Journal of global health. 2020. PMID: 32257154.
\bibitem{ref29} Jeong J. Effectiveness of the Moments that Matter parenting programme on caregiving practices, family violence, mental health and early child development in Kenya: a cluster-randomised controlled trial. BMJ global health. 2026. PMID: 42556995.
\bibitem{ref30} Ullman C. Interventions to prevent violence against women and girls globally: a global systematic review of reviews to update the RESPECT women framework. BMJ Public Health. 2025. PMID: 40017928.
\bibitem{ref31} Falb KL. Adapting, Sustaining and Scaling Violence against Women and Children Prevention Interventions in Low- and Middle-Income Countries: A Scoping Review on the Use of Implementation Science. Trauma, violence \& abuse. 2026. PMID: 40462272.
\bibitem{ref32} Doe Patience Fakornam. Couple-based preconception care as a gender-responsive strategy for improving health and well-being: a systematic review. Reproductive Health. 2026. PMID: 42316254.
\bibitem{ref33} Hawkins AJ. Exploring programmatic moderators of the effectiveness of marriage and relationship education programs: a meta-analytic study. Behavior Therapy. 2012. PMID: 22304880.
\bibitem{ref34} Bahrami-Samani S. Comparing the effects of premarital booklet- and video-based educations on the reproductive health literacy of engaged couples. Health Science Reports. 2024. PMID: 39377020.
\bibitem{ref35} Ashrafi E. Evaluating the Impact of Reproductive and Sexual Health Education on Knowledge and Attitudes of Couples in Premarital Counseling Classes: A Quasi-Experimental Study. Health Science Reports. 2026. PMID: 41773204.
\bibitem{ref36} Spencer CM. Online relationship education programs improve individual and relationship functioning: A meta-analytic review. Journal of Marital and Family Therapy. 2021. PMID: 33502054.
\bibitem{ref37} Robles TF. Marital quality and health: a meta-analytic review. Psychological Bulletin. 2014. PMID: 23527470.
\bibitem{ref38} O'Doherty L. Screening women for intimate partner violence in healthcare settings. Cochrane Database of Systematic Reviews. 2015. PMID: 26200817.
\bibitem{ref39} Feltner Cynthia. Screening for Intimate Partner Violence and for Caregiver Abuse of Older or Vulnerable Adults: An Evidence Report and Systematic Review for the US Preventive Services Task Force. JAMA. 2025. PMID: 40553460.
\bibitem{ref40} US Preventive Services Task Force. Screening for Intimate Partner Violence and Caregiver Abuse of Older or Vulnerable Adults: US Preventive Services Task Force Recommendation Statement. JAMA. 2025. PMID: 40553450.
\bibitem{ref41} Kalra Naira. Training healthcare providers to respond to intimate partner violence against women. Cochrane Database of Systematic Reviews. 2021. PMID: 34057734.
\bibitem{ref42} Betron. Effectiveness of a clinic-based counselling intervention on risk of experiencing intimate partner violence and reproductive coercion: a matched-pair cluster-controlled trial in Ebonyi and Sokoto, Nigeria. BMJ Global Health. 2025. PMID: 41151831.
\bibitem{ref43} World Health Organization. Responding to intimate partner violence and sexual violence against women: WHO clinical and policy guidelines. Geneva: World Health Organization; 2013. ISBN 978 92 4 154859 5. Available from: https://iris.who.int/handle/10665/85240; not PubMed-indexed.
\bibitem{ref44} World Health Organization. Health care for women subjected to intimate partner violence or sexual violence: a clinical handbook. Geneva: World Health Organization; 2014. WHO reference WHO/RHR/14.26. Available from: https://iris.who.int/handle/10665/136101; not PubMed-indexed.
\bibitem{ref45} Karakurt. Couples Therapy for Intimate Partner Violence: A Systematic Review and Meta-Analysis. Journal of marital and family therapy. 2016. PMID: 27377617.
\bibitem{ref46} McCollum. Couples treatment for interpersonal violence: a review of outcome research literature and current clinical practices. Violence and victims. 2008. PMID: 18624089.
\bibitem{ref47} Schneider. From scared to repaired: using an attachment-based perspective to understand situational couple violence. Journal of marital and family therapy. 2014. PMID: 25039387.
\bibitem{ref48} Chang. Health care interventions for intimate partner violence: what women want. Women's health issues. 2005. PMID: 15661584.
\bibitem{ref49} Robertson. Women and men's use of coercive control in intimate partner violence. Violence and victims. 2011. PMID: 21780535.
\bibitem{ref50} Swan. Factor structures for aggression and victimization among women who used aggression against male partners. Violence against women. 2012. PMID: 23012348.
\bibitem{ref51} Evans. "Even 'Daily' is Not Enough": How Well Do We Measure Domestic Violence and Abuse?-A Think-Aloud Study of a Commonly Used Self-Report Scale. Violence and victims. 2016. PMID: 26645540.
\bibitem{ref52} Skivington K. A new framework for developing and evaluating complex interventions: update of Medical Research Council guidance. BMJ. 2021. PMID: 34593508.
\bibitem{ref53} Moore G. Adapting interventions to new contexts-the ADAPT guidance. BMJ. 2021. PMID: 34344699.
\bibitem{ref54} Moore GF. Process evaluation of complex interventions: Medical Research Council guidance. BMJ. 2015. PMID: 25791983.
\bibitem{ref55} Horvat L. Cultural competence education for health professionals. Cochrane Database of Systematic Reviews. 2014. PMID: 24793445.
\bibitem{ref56} Vella Elizabeth. Does cultural competence training for health professionals impact culturally and linguistically diverse patient outcomes? A systematic review of the literature. Nurse Education Today. 2022. PMID: 35964378.
\bibitem{ref57} Hardy Billie-Jo. Systematic review of Indigenous cultural safety training interventions for healthcare professionals in Australia, Canada, New Zealand and the United States. BMJ Open. 2023. PMID: 37793931.
\bibitem{ref58} Tervalon M. Cultural humility versus cultural competence: a critical distinction in defining physician training outcomes in multicultural education. Journal of Health Care for the Poor and Underserved. 1998. PMID: 10073197.
\bibitem{ref59} Bresnahan Mary. Culturally safe healthcare: changing the lens from provider control to patient agency. Journal of Communication in Healthcare. 2024. PMID: 38426444.
\bibitem{ref60} Cano-Ibáñez Naomi. Physician-Patient Language Discordance and Poor Health Outcomes: A Systematic Scoping Review. Frontiers in Public Health. 2021. PMID: 33816420.
\bibitem{ref61} Chu Janet N. Association between language discordance and unplanned hospital readmissions or emergency department revisits: a systematic review and meta-analysis. BMJ Quality \& Safety. 2024. PMID: 38160059.
\bibitem{ref62} Heath Morten. Interpreter services and effect on healthcare - a systematic review of the impact of different types of interpreters on patient outcome. Journal of Migration and Health. 2023. PMID: 36816444.
\bibitem{ref63} Lauridsen Iben Gad. A systematic review of whether the number of linguistic errors in medical interpretation is associated with the use of professional vs ad hoc interpreters. Archives of Public Health. 2024. PMID: 39696623.
\bibitem{ref64} Reaume M. Patient-Physician Language Concordance and Cardiovascular Outcomes Among Patients With Hypertension. JAMA Network Open. 2025. PMID: 39969882.
\bibitem{ref65} Kosiyaporn Hathairat. Strengthening the migrant-friendliness of Thai health services through interpretation and cultural mediation: a system analysis. Global Health Research and Policy. 2020. PMID: 33372646.
\bibitem{ref66} König Andrea. A Systematic Scoping Review on Migrant Health Coverage in Thailand. Tropical Medicine and Infectious Disease. 2022. PMID: 36006258.
\bibitem{ref67} Phanwichatkul Titaree. The perceptions and practices of Thai health professionals providing maternity care for migrant Burmese women: An ethnographic study. Women and Birth. 2022. PMID: 34272187.
\bibitem{ref68} Buathong Napakkawat. Exploring symptoms perception and barriers to medication adherence among Thai Muslim patients with non-communicable diseases in a rural community in southern Thailand: a mixed-methods study. BMJ Open. 2024. PMID: 39653570.
\bibitem{ref69} Lohmae Userow. Muslim mothers' experiences in taking care of children with open heart surgery: A qualitative study in Southern Thailand. Belitung Nursing Journal. 2025. PMID: 40256382.
\bibitem{ref70} Pérez-Sánchez María. Access of migrant women to sexual and reproductive health services: A systematic review. Midwifery. 2024. PMID: 39243595.
\bibitem{ref71} Weschasat Tassanapan. The edutainment program on knowledge, perception, and uptake of cervical cancer screening among Muslim women in Southern Thailand: a quasi experimental study. BMC Public Health. 2024. PMID: 38971727.
\bibitem{ref72} Chuemchit Montakarn. Prevalence of Intimate Partner Violence in Thailand. Journal of Family Violence. 2018. PMID: 29904232.
\bibitem{ref73} Thananowan Nanthana. Intimate partner violence and factors associated with sexually transmitted infections among Thai women attending gynecology clinics. International Journal of STD \& AIDS. 2021. PMID: 33308089.
\bibitem{ref74} Thananowan Nanthana. Intimate partner violence among pregnant Thai women. Violence Against Women. 2008. PMID: 18408170.
\bibitem{ref75} Thitiyanviroj Benjaporn. Perceptions of Screening Women for Intimate Partner Violence Among Health Care Providers in Thailand. Nursing for Women's Health. 2024. PMID: 39536796.
\bibitem{ref76} Hayee Fusiyah. Sexual risk behaviors and influencing factors among Muslim adolescents on southern border of Thailand. International Journal of Adolescent Medicine and Health. 2020. PMID: 32549162.
\bibitem{ref77} Singkun Awirut. Risk of HIV/STIs among Muslim army conscripts in the three deep southern provinces of Thailand. Health Promotion Perspectives. 2021. PMID: 35079589.
\bibitem{ref78} Singkun Awirut. Sexual knowledge based on Islamic values and sexual risk behaviors of HIV/STIs among Thai Muslim army conscripts: A cross-sectional study. Belitung Nursing Journal. 2022. PMID: 37554485.
\bibitem{ref79} Ford Kathleen. Mental health in a conflict area: Migration, economic stress and religiosity in the three southernmost provinces of Thailand. International Journal of Social Psychiatry. 2017. PMID: 28024446.
\bibitem{ref80} Ford Kathleen. Mental health among youth in a conflict area: a cross-sectional study of the effects of parental absence, social media, family support and religious practice on psychiatric symptoms in Thailand's southernmost provinces. Conflict and Health. 2026. PMID: 41540416.
\bibitem{ref81} Vapattanawong Patama. A cross-sectional study on factors influencing adolescent depression in the civil-conflict areas of Thailand. BMC Public Health. 2026. PMID: 41872869.
\bibitem{ref82} Ford Kathleen. Coming of age in a conflict area: Mental health, education, employment, migration and family formation in the southernmost provinces of Thailand. International Journal of Social Psychiatry. 2018. PMID: 29417854.
\bibitem{ref83} Sirithammaphan Uraiwan. Barriers to measles mumps rubella vaccine acceptance in the three southern border provinces of Thailand. Clinical and Experimental Vaccine Research. 2023. PMID: 38025912.
\bibitem{ref84} Jinarong Taqwa. Muslim parents' beliefs and factors influencing complete immunization of children aged 0-5 years in a Thai rural community: a qualitative study. BMC Public Health. 2023. PMID: 37442996.
\bibitem{ref85} Napier AD. Culture and health. Lancet. 2014. PMID: 25443490.
\bibitem{ref86} Devakumar D. Racism, xenophobia, discrimination, and the determination of health. Lancet. 2022. PMID: 36502848.
\bibitem{ref87} Selvarajah S. Racism, xenophobia, and discrimination: mapping pathways to health outcomes. Lancet. 2022. PMID: 36502849.
\bibitem{ref88} Abubakar I. The UCL-Lancet Commission on Migration and Health: review of the state of progress. Lancet. 2026. PMID: 42314718.
\bibitem{ref89} Li S. Efficacy of culturally adapted interventions for common mental disorders in people of Chinese descent: a systematic review and meta-analysis. Lancet Psychiatry. 2023. PMID: 37208113.
\bibitem{ref90} Axinn WG. Community-Level Social Support Infrastructure and Adult Onset of Major Depressive Disorder in a South Asian Postconflict Setting. JAMA Psychiatry. 2022. PMID: 35080609.
\bibitem{ref91} Watt RG. Effectiveness and cost-effectiveness of a parenting programme to improve family wellbeing in England (TOGETHER): a multicentre, single-blind, randomised controlled trial. Lancet Public Health. 2026. PMID: 41887826.
\bibitem{ref92} Taft. Examining strength at home couples to prevent intimate partner violence on a military installation: A randomized controlled trial. Journal of consulting and clinical psychology. 2024. PMID: 38206858.
\bibitem{ref93} Ghuman SJ. Women's autonomy and child survival: a comparison of Muslims and non-Muslims in four Asian countries. Demography. 2003. PMID: 12962056.
\bibitem{ref94} Chen. Intimate partner violence: counseling, community resources, and legal issues for IPV victims and perpetrators. FP essentials. 2013. PMID: 24053261.
\bibitem{ref95} Decker. Implementing Trauma-Informed Partner Violence Assessment in Family Planning Clinics. Journal of women's health. 2017. PMID: 28375750.
\bibitem{ref96} MacMillan. Approaches to screening for intimate partner violence in health care settings: a randomized trial. JAMA. 2006. PMID: 16882959.
\bibitem{ref97} Ford-Gilboe. Development of a brief measure of intimate partner violence experiences: the Composite Abuse Scale (Revised)-Short Form (CASR-SF). BMJ open. 2016. PMID: 27927659.
\bibitem{ref98} Hegarty. Validation of the Coercive-Composite Abuse Scale (C-CAS) in an Australian sample of women experiencing intimate partner abuse. BMJ open. 2026. PMID: 42323273.
\bibitem{ref99} Wilson. Psychometric Properties of the Coercion in Intimate Partner Relationships Scale. Assessment. 2023. PMID: 34167331.
\bibitem{ref100} Sarac A. Cultural adaptation and pilot evaluation of the Couple Learning Program (CLP) among young Turkish couples: A mixed-method study. Acta Psychologica. 2026. PMID: 42442188.
\bibitem{ref101} Welland C. Culturally specific treatment for partner-abusive Latino men: a qualitative study to identify and implement program components. Violence and Victims. 2010. PMID: 21287968.
\bibitem{ref102} Sikoki B. A qualitative study on perceptions of adolescents' sexual and reproductive health education in Yogyakarta, Indonesia. International Journal of Adolescent Medicine and Health. 2024. PMID: 39395200.
\bibitem{ref103} Fauk NK. Cultural and religious determinants of HIV transmission: A qualitative study with people living with HIV in Belu and Yogyakarta, Indonesia. PLoS ONE. 2021. PMID: 34780506.
\bibitem{ref104} Grisurapong S. Establishing a one-stop crisis center for women suffering violence in Khonkaen hospital, Thailand. International Journal of Gynaecology and Obstetrics. 2002. PMID: 12429436.
\bibitem{ref105} Nagae M. [Domestic violence as a women's health issue--activities of health professionals in Thailand]. Japanese Journal of Public Health (Nihon Koshu Eisei Zasshi). 2004. PMID: 15162975.
\bibitem{ref106} Williamson HC. Does premarital education decrease or increase couples' later help-seeking? Journal of Family Psychology. 2014. PMID: 24294929.
\bibitem{ref107} Williamson HC. Barriers and facilitators of relationship help-seeking among low-income couples. Journal of Family Psychology. 2019. PMID: 30489129.
\bibitem{ref108} Oduenyi. Integrating gender-based violence first-line response into family planning and antenatal care services: an assessment of feasibility, appropriateness and acceptability to healthcare providers in Sokoto and Ebonyi States, Nigeria. BMC Health Services Research. 2025. PMID: 41267027.
\bibitem{ref109} Rhodes. The anatomy of a community health center system-level intervention for intimate partner violence. Journal of Urban Health. 2014. PMID: 23917943.
\bibitem{ref110} Miller Elizabeth. A family planning clinic-based intervention to address reproductive coercion: a cluster randomized controlled trial. Contraception. 2016. PMID: 26892333.
\bibitem{ref111} Zaher. Effect of domestic violence training: systematic review of randomized controlled trials. Canadian Family Physician. 2014. PMID: 25022633.
\bibitem{ref112} Jacubowski N. Marriage is not a safe place: heterosexual marriage and HIV-related vulnerability in Indonesia. Culture, Health \& Sexuality. 2008. PMID: 18038283.
\bibitem{ref113} Michau L. SASA! Together: An evolution of the SASA! approach to prevent violence against women. Evaluation and program planning. 2021. PMID: 33578229.
\bibitem{ref114} Stern E. Lessons learned from implementing Indashyikirwa in Rwanda - an adaptation of the SASA! approach to prevent and respond to intimate partner violence. Evaluation and program planning. 2018. PMID: 30125773.
\bibitem{ref115} Clyde TL. Revising Premarital Relationship Interventions for the Next Generation. Journal of Marital and Family Therapy. 2020. PMID: 30725473.
\bibitem{ref116} Silliman B. Improving practice in marriage preparation. Journal of Sex \& Marital Therapy. 1999. PMID: 10081741.
\bibitem{ref117} Paiboonsukwong Kittiphong. Thalassemia in Thailand. Hemoglobin. 2022. PMID: 35950590.
\bibitem{ref118} Zust BL. 10-Year Study of Christian Church Support for Domestic Violence Victims: 2005-2015. Journal of Interpersonal Violence. 2021. PMID: 33729071.
\bibitem{ref119} Hegarty K et al. Protocol for a randomised controlled trial of a web-based healthy relationship tool and safety decision aid for women experiencing domestic violence (I-DECIDE). BMC Public Health. 2015; PMID: 26231225.
\bibitem{ref120} Hegarty K et al. An online healthy relationship tool and safety decision aid for women experiencing intimate partner violence (I-DECIDE): a randomised controlled trial. The Lancet Public Health. 2019; PMID: 31155223.
\bibitem{ref121} Tarzia L et al. I-DECIDE: An Online Intervention Drawing on the Psychosocial Readiness Model for Women Experiencing Domestic Violence. Women's Health Issues. 2016; PMID: 26362841.
\bibitem{ref122} Tarzia L et al. Methodological and Ethical Challenges in a Web-Based Randomized Controlled Trial of a Domestic Violence Intervention. Journal of Medical Internet Research. 2017; PMID: 28351830.
\bibitem{ref123} Glass N et al. A safety app to respond to dating violence for college women and their friends: the MyPlan study randomized controlled trial protocol. BMC Public Health. 2015; PMID: 26350482.
\bibitem{ref124} Glass NE et al. Longitudinal Impact of the myPlan App on Health and Safety Among College Women Experiencing Partner Violence. Journal of Interpersonal Violence. 2022; PMID: 33576291.
\bibitem{ref125} Bloom TL et al. Concerned friends of intimate partner violence survivors: results from the myPlan randomized controlled trial on college campuses. BMC Public Health. 2023; PMID: 37259087.
\bibitem{ref126} Glass N et al. Effectiveness of the myPlan Teen App, a Digital Healthy Relationship and Safety Planning Intervention With Adolescent Aged 15-17 Years. Journal of Adolescent Health. 2024; PMID: 39066749.
\bibitem{ref127} Decker MR et al. Adapting the myPlan safety app to respond to intimate partner violence for women in low and middle income country settings: app tailoring and randomized controlled trial protocol. BMC Public Health. 2020; PMID: 32471469.
\bibitem{ref128} Decker MR et al. Safety decision-making and planning mobile app for intimate partner violence prevention and response: randomised controlled trial in Kenya. BMJ Global Health. 2020; PMID: 32675229.
\bibitem{ref129} Ford-Gilboe M et al. A tailored online safety and health intervention for women experiencing intimate partner violence: the iCAN Plan 4 Safety randomized controlled trial protocol. BMC Public Health. 2017; PMID: 28327116.
\bibitem{ref130} Ford-Gilboe M et al. Longitudinal impacts of an online safety and health intervention for women experiencing intimate partner violence: randomized controlled trial. BMC Public Health. 2020; PMID: 32098633.
\bibitem{ref131} O'Campo P et al. Design and Development of a Suite of Intimate Partner Violence Screening and Safety Planning Web Apps: User-Centered Approach. Journal of Medical Internet Research. 2021; PMID: 34931998.
\bibitem{ref132} Chuemchit M et al. Development and Preliminary Evaluation of iCanPlan: A Mobile Health Application for Intimate Partner Violence Prevention in Thailand. International Journal of Environmental Research and Public Health. 2026; PMID: 42196762.
\bibitem{ref133} Berry-Cabán CS et al. Empowering Safety: The Case for a Decision-making App for Military Partners Facing Intimate Partner Violence. Military Medicine. 2026; PMID: 42560214.
\bibitem{ref134} Ahmad F et al. Computer-assisted screening for intimate partner violence and control: a randomized trial. Annals of Internal Medicine. 2009; PMID: 19487706.
\bibitem{ref135} Lenert L et al. Electronic Health Record-Based Screening for Intimate Partner Violence: A Cluster Randomized Clinical Trial. JAMA Network Open. 2024; PMID: 39088215.
\bibitem{ref136} Wood SN et al. Risk for Severe Intimate Partner Violence in Nairobi's Informal Settlements: Tailoring the Danger Assessment to Kenya. Global Health, Science and Practice. 2024; PMID: 38290753.
\bibitem{ref137} El Morr C et al. Effectiveness of ICT-based intimate partner violence interventions: a systematic review. BMC Public Health. 2020; PMID: 32894115.
\bibitem{ref138} Anderson EJ et al. Web-Based and mHealth Interventions for Intimate Partner Violence Victimization Prevention: A Systematic Review. Trauma, Violence, \& Abuse. 2021; PMID: 31742475.
\bibitem{ref139} Ali K et al. Online Peer-to-Peer Support for Young People With Mental Health Problems: A Systematic Review. JMIR Mental Health. 2015; PMID: 26543923.
\bibitem{ref140} Marshall P et al. Understanding the Impacts of Online Mental Health Peer Support Forums: Realist Synthesis. JMIR Mental Health. 2024; PMID: 38722680.
\bibitem{ref141} Marshall P et al. Understanding Safety in Online Mental Health Forums: Realist Evaluation. JMIR Mental Health. 2025; PMID: 40577699.
\bibitem{ref142} Perry A et al. Suicidal behaviours and moderator support in online health communities: a scoping review. BMJ Open. 2021; PMID: 34193497.
\bibitem{ref143} Perowne R et al. Barriers and enablers to the moderation of self-harm content for a young person's online forum. Journal of Mental Health. 2024; PMID: 35574666.
\bibitem{ref144} Deng D et al. The Role of Moderators in Facilitating and Encouraging Peer-to-Peer Support in an Online Mental Health Community: A Qualitative Exploratory Study. Journal of Technology in Behavioral Science. 2023; PMID: 36810998.
\bibitem{ref145} Lerman K et al. Safe spaces or toxic places? Content moderation and social dynamics of online eating disorder communities. EPJ Data Science. 2025; PMID: 40726606.
\bibitem{ref146} Bender JL et al. How Cancer Online Support Groups Work, for Whom, and in What Circumstances: Realist Review. Journal of Medical Internet Research. 2026; PMID: 42127226.
\bibitem{ref147} Chu TH et al. Online Social Support for Intimate Partner Violence Victims in China: Quantitative and Automatic Content Analysis. Violence Against Women. 2021; PMID: 32339059.
\bibitem{ref148} Rempel E et al. Intimate partner violence: a review of online interventions. Informatics for Health \& Social Care. 2019; PMID: 29537928.
\bibitem{ref149} Rogers MM et al. Technology-Facilitated Abuse in Intimate Relationships: A Scoping Review. Trauma, Violence \& Abuse. 2023; PMID: 35537445.
\bibitem{ref150} Henry N et al. Technology-Facilitated Domestic and Sexual Violence: A Review. Violence Against Women. 2020; PMID: 32998673.
\bibitem{ref151} Henry N et al. Technology-Facilitated Domestic Violence Against Immigrant and Refugee Women: A Qualitative Study. Journal of Interpersonal Violence. 2022; PMID: 33719681.
\bibitem{ref152} Weekes CJ et al. Cyberstalking Perpetrators and Their Methods: A Systematic Literature Review. Trauma, Violence \& Abuse. 2026; PMID: 40271872.
\bibitem{ref153} Irving M et al. Post-separation Parenting Apps as Tools for Control and Resistance. Violence Against Women. 2026; PMID: 41823642.
\bibitem{ref154} Schlosser AE et al. Self-disclosure versus self-presentation on social media. Current Opinion in Psychology. 2020; PMID: 31357096.

\bibitem{ref156} Pretorius C et al. Young People's Online Help-Seeking and Mental Health Difficulties: Systematic Narrative Review. Journal of Medical Internet Research. 2019; PMID: 31742562.
\bibitem{ref157} Douglass CH et al. Social Media and Online Digital Technology Use Among Muslim Young People and Parents: Qualitative Focus Group Study. JMIR Pediatrics and Parenting. 2022; PMID: 35536616.
\bibitem{ref158} Robbins KC et al. The Independent Living Donor Advocate: An Essential Role for Living Kidney Donation. Nephrology Nursing Journal. 2014; PMID: 26287054.
\bibitem{ref159} Rudow DL et al. The Psychosocial and Independent Living Donor Advocate Evaluation and Post-surgery Care of Living Donors. Journal of Clinical Psychology in Medical Settings. 2015; PMID: 26293351.
\bibitem{ref160} Schenck D et al. Ethical Analysis and Policy Recommendations Regarding Domino Liver Transplantation. Transplantation. 2018; PMID: 29708521.
\bibitem{ref161} Sethi N et al. Ethics and the facial plastic surgeon. European Archives of Oto-Rhino-Laryngology. 2016; PMID: 26254909.
\bibitem{ref162} Pope KS et al. Dual relationships in psychotherapy. Ethics \& Behavior. 1991; PMID: 11649348.
\bibitem{ref163} Scopelliti J et al. Dual relationships in mental health practice: issues for clinicians in rural settings. Australian and New Zealand Journal of Psychiatry. 2004; PMID: 15555031.
\bibitem{ref164} Messing JT et al. The average predictive validity of intimate partner violence risk assessment instruments. Journal of Interpersonal Violence. 2013; PMID: 23262817.
\bibitem{ref165} Wang PL et al. Assessing the Danger: Validation of Taiwan Intimate Partner Violence Danger Assessment. Journal of Interpersonal Violence. 2015; PMID: 25315482.
\bibitem{ref166} Messing JT et al. Validation and adaptation of the danger assessment-5: A brief intimate partner violence risk assessment. Journal of Advanced Nursing. 2017; PMID: 28921610.
\bibitem{ref167} Messing JT et al. Accounting for Multiple Nonfatal Strangulation in Intimate Partner Violence Risk Assessment. Journal of Interpersonal Violence. 2022; PMID: 33280504.
\bibitem{ref168} Gerth J et al. [The Ontario Domestic Assault Risk Assessment (ODARA) - Validity und authorised German translation of an intimate partner violence screening tool]. Fortschritte der Neurologie-Psychiatrie. 2014; PMID: 25383928.
\bibitem{ref169} Jolliffe Simpson AD et al. Unpacking Multiagency Structured Professional Judgment Risk Assessments for Family Violence. Journal of Interpersonal Violence. 2023; PMID: 36710516.
\bibitem{ref170} Belfrage H et al. Assessment and management of risk for intimate partner violence by police officers using the Spousal Assault Risk Assessment Guide. Law and Human Behavior. 2012; PMID: 22471386.
\bibitem{ref171} Hogan NR et al. Investigating racial disparities in violence risk assessment using the Spousal Assault Risk Assessment Guide-Version 3. Psychological Assessment. 2024; PMID: 38512165.
\bibitem{ref172} Sanderson D et al. "Increasing Warm Handoffs: Optimizing Community Based Referrals in Primary Care Using QI Methodology. Journal of Primary Care \& Community Health. 2021; PMID: 34109884.
\bibitem{ref173} Loo S et al. Evaluating a social risk screening and referral program in an urban safety-net hospital emergency department. Journal of the American College of Emergency Physicians Open. 2023; PMID: 36704207.
\bibitem{ref174} Stevens GJ et al. Efficacy of a short message service brief contact intervention (SMS-SOS) in reducing repetition of hospital-treated self-harm. British Journal of Psychiatry. 2024; PMID: 38083861.
\bibitem{ref175} Radin AK et al. Comparative effectiveness of safety planning intervention with instrumental support calls (ISC) versus caring contacts (CC) two-way text messages: the SPARC Trial protocol. Contemporary Clinical Trials. 2023; PMID: 37321352.
\bibitem{ref176} Radin AK et al. Comparative effectiveness of two versions of a caring contacts intervention in healthcare providers, staff and patients. Journal of Affective Disorders. 2023; PMID: 36963515.

\bibitem{ref179} Chambers RA et al. Implementation Fidelity and Theory-Informed Dose Effects of a Teen Pregnancy Prevention Program for Native American Youth. Prevention Science. 2023; PMID: 37191932.

\bibitem{ref182} Roozenbeek J et al. Psychological inoculation improves resilience against misinformation on social media. Science advances. 2022; PMID: 36001675.
\bibitem{ref183} Appel RE et al. Psychological inoculation improves resilience to and reduces willingness to share vaccine misinformation. Scientific reports. 2025; PMID: 40820144.
\bibitem{ref184} Biddlestone M et al. Video inoculation against election misinformation across 12 EU nations. Communications psychology. 2026; PMID: 41845063.
\bibitem{ref185} Basol M et al. Good News about Bad News: Gamified Inoculation Boosts Confidence and Cognitive Immunity Against Fake News. Journal of cognition. 2020; PMID: 31934684.
\bibitem{ref186} Roozenbeek J et al. Technique-based inoculation against real-world misinformation. Royal Society open science. 2022; PMID: 35600423.
\bibitem{ref187} Roozenbeek J et al. Disentangling Item and Testing Effects in Inoculation Research on Online Misinformation: Solomon Revisited. Educational and psychological measurement. 2021; PMID: 37929261.
\bibitem{ref188} Kang Y et al. Educational and skills-based interventions for preventing relationship and dating violence in adolescents and young adults. The Cochrane database of systematic reviews. 2026; PMID: 42535453.
\bibitem{ref189} Kettrey HH et al. Can Attitudes Serve as Proxies for Behavioral Outcomes of Dating Violence Prevention Programs? Broader Lessons From a Pilot Evaluation of the Relationship Education Project. Violence and victims. 2023; PMID: 37011949.
\bibitem{ref190} Miller E et al. One-year follow-up of a coach-delivered dating violence prevention program: a cluster randomized controlled trial. American journal of preventive medicine. 2013; PMID: 23790995.
\bibitem{ref191} Abebe KZ et al. A cluster-randomized trial of a middle school gender violence prevention program: Design, rationale, and sample characteristics. Contemporary clinical trials. 2017; PMID: 28821469.
\bibitem{ref192} Cherrier C et al. [Supporting the applicability and transferability of the "Sortir Ensemble \& Se Respecter" program in France]. Sante publique (Vandoeuvre-les-Nancy, France). 2024; PMID: 38580464.
\bibitem{ref193} Bannon RS et al. The Bystander Approach to Sexual Assault Risk Reduction: Effects on Risk Recognition, Perceived Self-Efficacy, and Protective Behavior. Violence and victims. 2017; PMID: 28234197.
\bibitem{ref194} Peeren S et al. "Counteract the gaslighting" - a thematic analysis of open-ended responses about what women survivors of intimate partner sexual violence need from services. BMC women's health. 2024; PMID: 38336660.
\bibitem{ref195} Barnes M et al. Being silenced, loneliness and being heard: understanding pathways to intimate partner violence \& abuse in young adults. a mixed-methods study. BMC public health. 2022; PMID: 35974354.
\bibitem{ref196} Whitton SW et al. Disclosure and Help-Seeking Experiences of Sexual and Gender Minority Victims of Intimate Partner Violence: A Mixed-Methods Study. Journal of interpersonal violence. 2024; PMID: 37882155.
\bibitem{ref197} Brunton R et al. Further psychometric examination of the Intimate Partner Violence Internalized Stigma Scale (IPVIS; Brunton \& Harris, 2023): Examining criterion-related validity. Acta psychologica. 2024; PMID: 39536684.
\bibitem{ref198} Tian CY et al. Generic Health Literacy Measurements for Adults: A Scoping Review. International journal of environmental research and public health. 2020; PMID: 33114157.
\bibitem{ref199} Senahad N et al. Development of health literacy assessment tool for 9-10 Years old children in Thailand. Public health in practice (Oxford, England). 2023; PMID: 37727706.
\bibitem{ref200} Couture-Carron A et al. One Size Doesn't Fit All: Dating Abuse Against Women From the Perspective of South Asian Muslim Youth in Canada. Journal of interpersonal violence. 2017; PMID: 26289454.
\bibitem{ref201} Thongpriwan V et al. Reflections on attitudes, experiences, and vulnerability of intimate partner violence among Southeast Asian college women living in United States. Asian journal of psychiatry. 2015; PMID: 26442989.
\bibitem{ref202} Bell CA et al. Humility, Forgiveness, and Emerging Adult Female Romantic Relationships. Journal of marital and family therapy. 2019; PMID: 29073326.
\bibitem{ref203} Koetke J et al. Intellectual humility is reliably associated with constructive responses to conflict. PloS one. 2024; PMID: 39240981.
\bibitem{ref204} Lehmann M et al. Intellectual Humility Predicts Empathic Accuracy and Empathic Resilience. Personality \& social psychology bulletin. 2026; PMID: 39902736.
\bibitem{ref205} Jongman-Sereno KP et al. Intellectual Humility and Responsiveness to Public Health Recommendations. Personality and individual differences. 2023; PMID: 37426514.
\bibitem{ref206} Foronda C et al. A Theory of Cultural Humility. Journal of transcultural nursing. 2020; PMID: 31516091.
\bibitem{ref207} Lekas HM et al. Rethinking Cultural Competence: Shifting to Cultural Humility. Health services insights. 2020; PMID: 33424230.
\bibitem{ref208} Sotto-Santiago S et al. S.E.L.F.: Lessons in institutional cultural humility approach from an academic health center. Dialogues in health. 2026; PMID: 42656867.
\bibitem{ref209} MacBeth A et al. Exploring compassion: a meta-analysis of the association between self-compassion and psychopathology. Clinical psychology review. 2012; PMID: 22796446.
\bibitem{ref210} Millard LA et al. The effectiveness of compassion focused therapy with clinical populations: A systematic review and meta-analysis. Journal of affective disorders. 2023; PMID: 36649790.
\bibitem{ref211} Baker LR et al. Self-compassion and relationship maintenance: the moderating roles of conscientiousness and gender. Journal of personality and social psychology. 2011; PMID: 21280964.
\bibitem{ref212} Kaya F et al. The predictive effect of self-compassion on relationship satisfaction and conflict resolution styles in romantic relationships in nursing students. Nursing forum. 2022; PMID: 35245387.
\bibitem{ref213} Williams Z et al. Relationship satisfaction in Black heterosexual couples: The role of self-compassion and openness. Journal of marital and family therapy. 2023; PMID: 37752743.
\bibitem{ref214} Godfrey DA et al. Empathy Mediates the Relations between Working Memory and Perpetration of Intimate Partner Violence and Aggression. Behavioral sciences (Basel, Switzerland). 2020; PMID: 32150915.
\bibitem{ref215} Brassard A et al. Childhood Sexual Abuse, Dyadic Empathy, and Intimate Partner Violence Among Men Seeking Psychological Help. Journal of interpersonal violence. 2022; PMID: 35089108.
\bibitem{ref216} Buitelaar NJL et al. ADHD in Childhood and/or Adulthood as a Risk Factor for Domestic Violence or Intimate Partner Violence: A Systematic Review. Journal of attention disorders. 2020; PMID: 25995243.
\bibitem{ref217} Boonyananth N et al. A Systematic Review of the Psychosocial Mechanism Underlying the Pathway from Intra-Familial Victimization in Childhood to Intimate Relationship Violence in Adolescence and Adulthood. Trauma, violence \& abuse. 2026; PMID: 40035136.
\bibitem{ref218} Gopalakrishnan L et al. Psychometric Validation of Trust, Commitment, and Satisfaction Scales to Measure Marital Relationship Quality Among Newly Married Women in Nepal. International journal of environmental research and public health. 2025; PMID: 41007600.
\bibitem{ref219} Körün AB et al. Even Though the Long Distance: Are We Still Going on? Dyadic Trust, Relationship Maintenance Behaviors, and Relationship Quality Among Emerging Adulthoods. Scandinavian journal of psychology. 2026; PMID: 41235925.
\bibitem{ref220} Gonçalves JPB et al. The role of religiosity and spirituality in interpersonal violence: a systematic review and meta-analysis. Revista brasileira de psiquiatria (Sao Paulo, Brazil). 2023; PMID: 36331229.
\bibitem{ref221} Simonič B et al. The Power of Women's Faith in Coping with Intimate Partner Violence: Systematic Literature Review. Journal of religion and health. 2021; PMID: 33704630.

\bibitem{ref223} Ripley JS et al. Obtaining religious and spiritual competencies for relationship therapy: Outcomes of Competency Addressing Religion and Spirituality (CARS) training. Psychotherapy (Chicago, Ill.). 2026; PMID: 40705620.
\bibitem{ref224} Aziz AFA et al. Knowledge improvement on premarital HIV counseling with PAUSE flipchart versus conventional method: a cluster randomised controlled trial. BMC primary care. 2026; PMID: 42420861.
\bibitem{ref225} Mustafa Y et al. Islam and the four principles of medical ethics. Journal of Medical Ethics. 2014; PMID: 23975951.
\bibitem{ref226} Padela AI et al. Muslim patients and cross-gender interactions in medicine: an Islamic bioethical perspective. Journal of Medical Ethics. 2011; PMID: 21041237.
\bibitem{ref227} Doedes E et al. Moral universe of Muslim healthcare practitioners in the UK: balancing Islamic and secular ethics in palliative and end-of-life care. Journal of Medical Ethics. 2026; PMID: 41651650.
\bibitem{ref228} Padela AI et al. Islamic medical ethics: a primer. Bioethics. 2007; PMID: 17845488.
\bibitem{ref229} Alsolamy S et al. Islamic views on artificial nutrition and hydration in terminally ill patients. Bioethics. 2014; PMID: 22845721.
\bibitem{ref230} Daar AS et al. Bioethics for clinicians: 21. Islamic bioethics. CMAJ (Canadian Medical Association Journal). 2001; PMID: 11202669.
\bibitem{ref231} Bhowmik B et al. Faith-Based Lifestyle Intervention for Diabetes Prevention Among Adults in Bangladesh: A Cluster Randomized Clinical Trial. JAMA Network Open. 2025; PMID: 41114978.
\bibitem{ref232} Zoellner LA et al. Lay-Led Intervention for War and Refugee Trauma: A Randomized Clinical Trial. JAMA Network Open. 2024; PMID: 39186273. Two errata published: JAMA Netw Open. 2024;7(9):e2440198 and 2024;7(10):e2446257.
\bibitem{ref233} Gendler Y et al. Exploring Cultural and Religious Effects on HPV Vaccination Decision Making Using a Web-Based Decision Aid: A Quasi-experimental Study. Medical Decision Making. 2024; PMID: 38600776.
\bibitem{ref234} Al-Sayaghi KM et al. Stressors, Needs and Satisfaction of Families of Critically Ill Arabic Patients: A Systematic Review. Nursing in Critical Care. 2026; PMID: 41379069.
\bibitem{ref235} Mohamed Mahdi Z et al. Gender-Based Violence Against Migrant Women From Islamic Background: A Systematic Review of International Studies. Trauma, Violence \& Abuse. 2025; PMID: 41139992.
\bibitem{ref236} Chamsi-Pasha H et al. Assisted reproductive technology: Islamic Sunni perspective. Human Fertility. 2015; PMID: 25660098.
\bibitem{ref237} Pan SW et al. Health Effects of Religion, Spirituality, and Supernatural Beliefs in Mainland China: A Systematic Review. Journal of Religion and Health. 2022; PMID: 35347576.
\bibitem{ref238} Abdel-Khalek AM et al. The Relationship Between Religiosity and Anxiety: A Meta-analysis. Journal of Religion and Health. 2019; PMID: 31309442.
\bibitem{ref239} Cleary M et al. Drivers and consequences of self-immolation in parts of Iran, Iraq and Uzbekistan: A systematic review of qualitative evidence. Burns. 2021; PMID: 31928787.
\bibitem{ref240} Mittal SO et al. Faith, fasting, and well-being: Emirates Neurology Society consensus guidelines on safe Ramadan fasting in Parkinson's disease. Frontiers in Neurology. 2025; PMID: 41458120.
\bibitem{ref241} Packer S et al. Informed consent with a focus on islamic views. The Journal of IMA (Islamic Medical Association of North America) (professional-association journal). 2011; PMID: 23610513.
\bibitem{ref242} Arawi TA et al. The muslim physician and the ethics of medicine. The Journal of IMA (Islamic Medical Association of North America) (professional-association journal). 2010; PMID: 23864762.
\bibitem{ref243} Irfan B et al. Considering Islamic Frameworks to Infectious Disease Prevention. Open Forum Infectious Diseases. 2025; PMID: 41141460.
\bibitem{who_vaw_2026} World Health Organization. Violence against women [fact sheet]. 2026. \url{https://www.who.int/news-room/fact-sheets/detail/violence-against-women}
\bibitem{dwf2024} Department of Women's Affairs and Family Development, Thailand. Domestic violence [in Thai]. 2024. \url{https://www.dwf.go.th/contents/71825} (accessed 2026-09-02)
\bibitem{who_respect} World Health Organization. RESPECT women: preventing violence against women. 2nd ed. Geneva: WHO; 2025. ISBN 9789240117020. Not PubMed-indexed.
\bibitem{fawcett2010} [Not PubMed-indexed; cited for transparency only.] Fawcett EB, Hawkins AJ, Blanchard VL, Carroll JS. Do premarital education programs really work? A meta-analytic study. \emph{Fam Relat}. 2010;59(3):232--239. \url{https://doi.org/10.1111/j.1741-3729.2010.00598.x}

\bibitem{charoenwong2025} Charoenwong W. Islamic premarital education in Malaysia through the lens of moral education. \emph{Asia-Pacific Social Science Review}. 2025. (Volume and page numbers not yet assigned by the publisher at time of writing.)
\end{thebibliography}
